\chapter{The mental model}

GitHub Actions is built on a small set of concepts.
If you understand these, everything else becomes configuration.

\section{Events and workflows}
An \emph{event} happens in the repository: a push, a pull request, a tag, or a schedule.
A \emph{workflow} declares which events it listens to and what jobs to run.

\begin{patternbox}{Pattern: start from the event}
Describe the event first, then design the workflow around it.
When the event is clear, jobs and steps follow naturally.
\end{patternbox}

\section{Jobs and steps}
A workflow runs one or more \emph{jobs}.
Each job runs on its own runner and is made of sequential \emph{steps}.
Jobs run in parallel by default unless you declare dependencies with \code{needs}.

\begin{lstlisting}[language=YAML]
jobs:
  build:
    runs-on: ubuntu-latest
    steps:
      - run: ./scripts/build

  test:
    runs-on: ubuntu-latest
    needs: build
    steps:
      - run: ./scripts/test
\end{lstlisting}

\section{Runners}
A runner is the machine that executes your job.
GitHub-hosted runners are ephemeral; each job starts on a clean VM.
Self-hosted runners are long-lived and must be secured and maintained.

\section{Contexts and expressions}
Expressions read data from \emph{contexts} like \code{github}, \code{inputs}, or \code{secrets}.
They are written as \code{\$\{\{ ... \}\}} and evaluated by the workflow engine.

\section{Environment variables and outputs}
Steps can read and write environment variables, and can publish outputs.
Outputs are the cleanest way to pass data across steps or jobs.

\begin{lstlisting}[language=YAML]
steps:
  - name: Compute version
    id: meta
    run: echo "version=1.2.3" >> "$GITHUB_OUTPUT"

  - name: Use version
    run: echo "Using ${{ steps.meta.outputs.version }}"
\end{lstlisting}

\section{Artifacts}
Artifacts are files produced by a job that you want to keep.
They are stored by GitHub and can be downloaded later.
Use artifacts for test reports, build outputs, and deployment packages.

\section{Concurrency and cancellation}
Large repositories need guardrails.
Use concurrency groups to avoid running the same workflow twice for the same branch.

\begin{lstlisting}[language=YAML]
concurrency:
  group: ${{ github.workflow }}-${{ github.ref }}
  cancel-in-progress: true
\end{lstlisting}

\section{Putting it together}
Think of a workflow as a thin controller:
\begin{itemize}[nosep]
  \item inputs come from events and contexts
  \item steps call scripts or actions
  \item outputs are artifacts or job outputs
\end{itemize}

\section{State and isolation}
Each job starts on a clean machine.
Assume nothing persists between jobs unless you use artifacts or caches.
This is why moving logic into scripts and storing outputs explicitly matters.

\section{Event payload as input}
The event payload is just data.
Treat it like configuration that drives a workflow, not as business logic.
If you need branching behavior, prefer small conditional steps over large, hidden logic.

\begin{checklistbox}{Checklist: mental model}
\begin{itemize}[nosep]
  \item Events select workflows.
  \item Workflows run jobs.
  \item Jobs run steps on runners.
  \item Outputs connect steps and jobs.
  \item Artifacts persist files outside the runner.
\end{itemize}
\end{checklistbox}
